\documentclass[10pt]{article}
\usepackage[margin=0.82in]{geometry}
\usepackage[T1]{fontenc}
\usepackage{lmodern}
\usepackage{amsmath,amssymb,amsthm,mathtools}
\usepackage{enumitem}
\usepackage{hyperref}

\newtheorem{theorem}{Theorem}
\newtheorem{corollary}{Corollary}
\newcommand{\ellT}{\ell_2(\Lambda;\mathbb R^M)}
\newcommand{\blocksection}[1]{%
  \par\addvspace{2.2ex}\noindent{\Large\bfseries #1}\par\nobreak\vspace{0.9ex}}

\title{The Optimized Curvelet ``Length Term'' Vanishes:\\
A Counterexample to the Final-Remarks Conjecture}
\author{Run fa\_banach\_001, agent\_lane\_14}
\date{11 August 2026}

\begin{document}
\maketitle

\blocksection{Status and exact source conjecture}

\textbf{Status:} \texttt{candidate\_counterexample\_likely\_valid}.

Fornasier and Rauhut's 2006 paper \cite{FR} introduces coefficient vectors
$u=(u_\lambda)_\lambda$, adaptive nonnegative weights
$v=(v_\lambda)_\lambda$, and the term
\begin{equation}\label{eq:Q}
 Q_{\theta,\rho}(u,v)
 =\sum_{\lambda\in\Lambda}
 \theta_\lambda(\rho_\lambda-v_\lambda)^2.
\end{equation}
In the final remarks on PDF page~36, the paper takes
$\omega_\lambda=\varepsilon$ and
$\theta_\lambda=1/(4\varepsilon)$ and conjectures that \eqref{eq:Q}, after
the adaptive minimization, estimates the length of a rectifiable curved
discontinuity as $\varepsilon\downarrow0$.

The conjectured term actually vanishes for every uniformly
$\ell_2$-bounded coefficient family.  This follows from the source's exact
weight update and does not depend on the choice of $\rho_\lambda$.

\blocksection{Exact no-go theorem}

For $1\le q\le\infty$, set
$a_\lambda=\|u_\lambda\|_q$.  Equation~(12) of \cite{FR} states that the
minimizer in $v_\lambda\ge0$, with $u$ fixed, is
\begin{equation}\label{eq:v}
 v_\lambda^*
 =\max\left\{\rho_\lambda-
 \frac{a_\lambda}{2\theta_\lambda},0\right\}.
\end{equation}

\begin{theorem}[collapse after optimizing the weights]\label{thm:nogo}
Let $\Lambda$ be countable, $M<\infty$, $u\in\ellT$, and
$\rho_\lambda\ge0$.  If $v^*$ is given by \eqref{eq:v}, then
\begin{equation}\label{eq:generalbound}
 Q_{\theta,\rho}(u,v^*)
 \le \sum_{\lambda\in\Lambda}
 \frac{\|u_\lambda\|_q^2}{4\theta_\lambda}.
\end{equation}
In particular, if $\theta_\lambda=1/(4\varepsilon)$, then
\begin{equation}\label{eq:epsbound}
 Q_{\varepsilon}(u,v^*)
 \le C_{M,q}\varepsilon\|u\|_{\ellT}^2,
 \qquad
 C_{M,q}=M^{\max\{2/q-1,0\}},
\end{equation}
with $C_{M,\infty}=1$.  Consequently
$Q_\varepsilon(u^\varepsilon,v^{*,\varepsilon})\to0$ for every family
$u^\varepsilon$ bounded in $\ellT$.
\end{theorem}

\begin{proof}
Formula \eqref{eq:v} gives the exact identity
\[
 \rho_\lambda-v_\lambda^*
 =\min\left\{\rho_\lambda,
 \frac{a_\lambda}{2\theta_\lambda}\right\}.
\]
Therefore, term by term,
\[
 \theta_\lambda(\rho_\lambda-v_\lambda^*)^2
 \le \frac{a_\lambda^2}{4\theta_\lambda},
\]
and summation proves \eqref{eq:generalbound}.  Under the proposed scaling,
the right side is $\varepsilon\sum_\lambda\|u_\lambda\|_q^2$.
Finite-dimensional norm comparison gives
\[
 \|x\|_q^2\le
 M^{\max\{2/q-1,0\}}\|x\|_2^2,
\]
which proves \eqref{eq:epsbound} and the limit.
\end{proof}

The estimate is pointwise in $\lambda$.  It remains true for arbitrary
$\varepsilon$-dependent thresholds $\rho_\lambda$, and restricting the sum
to curvelets at scales $2^{-j}\lesssim\varepsilon$ or to curvelets tangent
to the edge only makes the left side smaller.

\blocksection{Concrete rectifiable-edge counterexample}

Let $D$ be a disk compactly contained in the image domain and let
$f=\mathbf1_D$ in one channel (with the other channels zero, if $M>1$).
Its discontinuity set is $\partial D$, a rectifiable curve with
\[
 \mathcal H^1(\partial D)=2\pi r>0.
\]
Let $u=(u_\lambda)_\lambda$ be its curvelet analysis coefficients.  Since a
curvelet system is a frame, its upper frame bound $B$ gives
\[
 \sum_\lambda |u_\lambda|^2\le B\|\mathbf1_D\|_{L^2}^2<\infty.
\]
For every $\varepsilon>0$, choose any nonnegative thresholds
$\rho_\lambda^\varepsilon$---including the scale laws suggested in the
source---and optimize $v$ by \eqref{eq:v}.  Theorem~\ref{thm:nogo} yields
\[
 0\le Q_\varepsilon(u,v^{*,\varepsilon})
 \le B\varepsilon\|\mathbf1_D\|_{L^2}^2
 \longrightarrow0,
\]
whereas the discontinuity length stays $2\pi r$.  Thus the proposed term
cannot converge to the length, nor to any fixed positive multiple of it.

More generally, the same conclusion applies to any stable recovery family
whose coefficient norms stay bounded and whose limiting image has positive
rectifiable jump length.  For example, uniform coefficient boundedness for
identity data follows immediately by comparing a global minimizer with the
competitor $u=0$, $v=\rho$.  If a recovery family must instead grow at least
like $\varepsilon^{-1/2}$ merely to avoid \eqref{eq:epsbound}, it has lost the
stability required of the proposed morphological estimator.

\blocksection{Upgrade audit}

The counterexample was tested against five possible repairs.
\begin{itemize}
\item Changing $\rho_\lambda$ cannot help because it has disappeared from
the upper bound.
\item Selecting only fine-scale, overlapping, or tangent curvelets cannot
help because this passes to a sub-sum.
\item Joint channels and all finite $q$ only alter the harmless constant
$C_{M,q}$.
\item Replacing a fixed signal by any uniformly bounded minimizer family
still forces the term to zero.
\item A nonzero perimeter limit would require a different coupling or an
additional renormalization after eliminating $v$; that is a different
functional, not a proof of the conjecture as stated.
\end{itemize}

Exact-phrase, title, and arXiv searches through 11 August 2026 found no later
resolution of this particular conjecture.  The proof above uses only the
source's equation~(12), finite-dimensional norm equivalence, and the curvelet
frame upper bound.  The accompanying finite-sequence verifier checks the
exact identity and both inequalities for multiple $M,q,\varepsilon$, and
random positive thresholds.

\begin{thebibliography}{9}
\bibitem{FR}
M. Fornasier and H. Rauhut,
\emph{Recovery algorithms for vector valued data with joint sparsity
constraints}, arXiv:math/0608124 (2006).
\end{thebibliography}

\end{document}
